\documentclass[11pt]{article}
\usepackage[a4paper, top=18mm, bottom=18mm, left=20mm, right=20mm]{geometry}
\usepackage[T2A]{fontenc}
\usepackage[utf8]{inputenc}
\usepackage[ukrainian]{babel}
\usepackage{graphicx}
\setlength{\parindent}{0pt}
\setlength{\parskip}{4pt}
\pagestyle{empty}
\begin{document}
%begin
{\large\textbf{Тест mixed}}

\textbf{Варіант \No\,1}\hfill \textit{ПІБ:}\,\underline{\hspace{60mm}}
\vspace{4mm}

%beginitem
1. Відмітити все, що є однією правильною лексемою:
( 1 ) 8,3

( 2 ) <>

( 3 ) 8.3

( 4 ) FIFO

%enditem
%beginitem
2. Відмітити все, що є коректним оператором С++:
( 1 ) cout>>3;

( 2 ) int x+y=z;

( 3 ) cout<<3.4;

( 4 ) int f(int n);

%enditem
%beginitem
3. Що буде виведено за виконання фрагменту коду?

%code
#include <iostream>
using namespace std;

int f(int c){
    int w = 72;
    if (c > -2) 
        return 2;
    if (c >= 4)
         w = 9;
    else
         w = 0;
    return w;
}

int main(){
    cout << f(-6);
    return 0;
}
%code


%enditem



%end
\vspace{10mm}
\noindent\rule{\textwidth}{0.3pt}
\vspace{8mm}
%begin
{\large\textbf{Тест mixed}}

\textbf{Варіант \No\,2}\hfill \textit{ПІБ:}\,\underline{\hspace{60mm}}
\vspace{4mm}

%beginitem
1. Відмітити все, що є коректним оператором С++:
( 1 ) int f(int a,b);

( 2 ) if a<b a=0;

( 3 ) x-=2;

( 4 ) int x='a'-'c';

%enditem
%beginitem
2. Відмітити все, що є однією правильною лексемою:
( 1 ) [1]

( 2 ) =<

( 3 ) z13

( 4 ) >

%enditem
%beginitem
3. Що буде виведено за виконання фрагменту коду?

%code
#include <iostream>
using namespace std;

int h(int d){
    int x = 31;
    if (d < 5) 
        return 5;
    else if (d != 1)
         x = 4;
    else
         return 6;
    return x;
}

int main(){
    cout << h(0);
    return 0;
}
%code


%enditem



%end
\newpage
%begin
{\large\textbf{Тест mixed}}

\textbf{Варіант \No\,3}\hfill \textit{ПІБ:}\,\underline{\hspace{60mm}}
\vspace{4mm}

%beginitem
1. Відмітити все, що є однією правильною лексемою:
( 1 ) **

( 2 ) (1)

( 3 ) /

( 4 ) "1.1.1"

%enditem
%beginitem
2. Що буде виведено за виконання фрагменту коду?

%code
#include <iostream>
using namespace std;

int g(int a){
    int u = 37;
    if (a) 
        u = 7;
    else if (a == -4)
         return 1;
    else
         return 3;
    return u;
}

int main(){
    cout << g(1);
    return 0;
}
%code


%enditem
%beginitem
3. Відмітити все, що є коректним оператором С++:
( 1 ) cin<<x;

( 2 ) int 3a=2;

( 3 ) x=int(4);

( 4 ) x+=y;

%enditem



%end
\vspace{10mm}
\noindent\rule{\textwidth}{0.3pt}
\vspace{8mm}
%begin
{\large\textbf{Тест mixed}}

\textbf{Варіант \No\,4}\hfill \textit{ПІБ:}\,\underline{\hspace{60mm}}
\vspace{4mm}

%beginitem
1. Відмітити все, що є однією правильною лексемою:
( 1 ) False

( 2 ) x*y

( 3 ) 'c'

( 4 ) +>

%enditem
%beginitem
2. Що буде виведено за виконання фрагменту коду?

%code
#include <iostream>
using namespace std;

int g(int b){
    int z = 95;
    if (b) 
        z = 8;
    if (b <= 0)
         return 6;
    else
         z = 4;
    return z;
}

int main(){
    cout << g(6);
    return 0;
}
%code


%enditem
%beginitem
3. Відмітити все, що є коректним оператором С++:
( 1 ) bool r=('N'>'5');

( 2 ) "a"=6;

( 3 ) cout<<"\n";

( 4 ) 'a'=3;

%enditem



%end
\newpage
\end{document}

